\par\noindent\textbf{(1)} 周期边界的波矢态间隔 \(2\pi/L\)，每 \(k\) 含两个自旋。填满 \(-k_F\le k\le k_F\) 后 \(N=2Lk_F/\pi\)，所以 \(k_F=\pi n/2\)（\(n=N/L\)），\(E_F=\hbar^2\pi^2n^2/(8m)\)。
\par\noindent\textbf{(2)} 每个正负 \(k\) 与双自旋共计 \[g(\varepsilon)/L=(2/\pi)(dk/d\varepsilon)=(2/\pi)(m/(\hbar^2k))=\frac{\sqrt{2m}}{\pi\hbar\sqrt\varepsilon}\]。态密度在低能阈值呈 \(\varepsilon^{-1/2}\) 发散，属于一维van Hove型奇异。
\par\noindent\textbf{(3)} 对一般简并费米气体，低温热容 \(C_V=(\pi^2/3)g(E_F)k_B^2T+O(T^3)\)，其中 \(g\) 包含全部自旋。代入 \(g(E_F)=N/(2E_F)\)，得 \(C_V=(\pi^2/6)Nk_B^2T/E_F\)，不同于三维电子气系数。
\par\noindent\textbf{(4)} 零温内能 \(U_0=2(L/\pi)\int_0^{k_F}\hbar^2k^2/(2m)dk=NE_F/3\propto N^3/L^2\)。固定 \(N\) 求压强的一维对应张力 \(P_{1D}=-(\partial U_0/\partial L)_N=2U_0/L\propto n^3\)。这与三维非相对论的 \(n^{5/3}\) 多方律不同。
